\def\piedfiche{TD n°1 · 1ère S2}
\documentclass[10pt]{article}
\usepackage{amsmath,amssymb}
\usepackage{fontspec}

\setmainfont{Latin Modern Roman}
\newfontfamily\titre{Inter}[UprightFont=* Bold,BoldFont=* ExtraBold]
\usepackage[a4paper,top=1.1cm,bottom=1.35cm,left=1.4cm,right=1.4cm,headheight=0pt,footskip=0.8cm]{geometry}
\usepackage{graphicx}
\usepackage[most]{tcolorbox}
\usepackage{tikz}
\usepackage{enumitem}
\usepackage{fancyhdr}
\usepackage{eso-pic}
\usepackage{multicol}
\usepackage{lastpage}\usepackage{needspace}

% ---------- couleurs Sigma Cap ----------
\definecolor{ink}{HTML}{14284B}
\definecolor{sky}{HTML}{1F8FD0}
\definecolor{lime}{HTML}{6FA82F}
\definecolor{gold}{HTML}{D99A1E}
\definecolor{rouge}{HTML}{C0392B}
\definecolor{violet}{HTML}{6C3FB5}
\definecolor{fond}{HTML}{F2F5F9}

\setlength{\parindent}{0pt}
\setlength{\parskip}{2pt}
\renewcommand{\arraystretch}{1.15}
\newcommand{\R}{\mathbb{R}}
\newcommand{\N}{\mathbb{N}}
\newcommand{\Z}{\mathbb{Z}}
\newcommand{\Q}{\mathbb{Q}}

% ---------- filigrane ----------
\AddToShipoutPictureBG{%
  \AtPageCenter{\makebox(0,0){\rotatebox{50}{\fontsize{78}{78}\selectfont\bfseries\color{ink!6}SIGMA CAP}}}%
  \AtPageLowerLeft{\put(\LenToUnit{0.5\paperwidth},\LenToUnit{0.28\paperheight}){\makebox(0,0){\rotatebox{50}{\fontsize{30}{30}\selectfont\bfseries\color{sky!8}PROFDEMATHS 2.0}}}}%
  \AtPageUpperLeft{\put(0,\LenToUnit{-0.35cm}){\color{ink}\rule{\paperwidth}{0.35cm}}}%
  \AtPageLowerLeft{\put(0,0){\color{lime}\rule{\paperwidth}{0.18cm}}}%
}

% ---------- pied de page ----------
\pagestyle{fancy}\fancyhf{}
\renewcommand{\headrulewidth}{0pt}
\fancyfoot[L]{\scriptsize\color{ink}\textbf{Sigma Cap} -- Vers la réussite · CEM Kennedy}
\fancyfoot[C]{\scriptsize\color{ink!70}\piedfiche}
\fancyfoot[R]{\scriptsize\color{ink}\textbf{\thepage/\pageref{LastPage}}}

% ---------- puces colorées ----------
\setlist[itemize]{leftmargin=1.1em,itemsep=1pt,topsep=2pt,label={\color{sky}\small$\blacktriangleright$}}
\DeclareRobustCommand{\pucenum}[1]{\tikz[baseline=(n.base)]\node[circle,fill=sky,text=white,inner sep=1.3pt,minimum size=1.15em,font=\footnotesize\bfseries](n){#1};}
\setlist[enumerate,1]{leftmargin=1.6em,itemsep=1pt,topsep=1pt,label={\pucenum{\arabic*}}}
\setlist[enumerate,2]{leftmargin=1.5em,itemsep=1pt,topsep=1pt,label={\color{gold}\bfseries\alph*)}}

% ---------- cadres ----------
\newcounter{exo}
\newtcolorbox{exo}[3][]{enhanced,breakable,colback=white,colframe=#2,boxrule=0.9pt,arc=2.5mm,
  left=7pt,right=7pt,top=5pt,bottom=4pt,before skip=3.5pt,after skip=3pt,
  fonttitle=\bfseries\small,coltitle=white,colbacktitle=#2,
  title={\stepcounter{exo}Exercice \theexo\ \ \textcolor{white!85!#2}{\textbullet}\ \ #3},
  attach boxed title to top left={xshift=8pt,yshift=-2.2mm},
  boxed title style={arc=1.5mm,boxrule=0pt},#1}

\newcommand{\partie}[3]{\needspace{7\baselineskip}\vspace{3pt}%
  \begin{tcolorbox}[enhanced,colback=#1,colframe=#1,arc=1.5mm,boxrule=0pt,left=6pt,right=6pt,top=2.5pt,bottom=2.5pt,before skip=4pt,after skip=3pt]
  {\color{white}\bfseries\large #2}\hfill{\color{white}\small\itshape #3}
  \end{tcolorbox}}

\newcommand{\demo}{\tikz[baseline=-0.6ex]\node[fill=violet,text=white,rounded corners=1.5pt,inner sep=1.6pt,font=\scriptsize\bfseries]{DÉMONSTRATION};}
\newcommand{\etoiles}[1]{{\color{gold}\ifcase#1\or$\star$\or$\star\star$\or$\star\star\star$\fi}}

% bulle avec la photo de M. Thiam
\newcommand{\bulle}[1]{%
  \begin{tcolorbox}[enhanced,colback=gold!9,colframe=gold!80!black,boxrule=0.6pt,arc=3mm,
    left=34pt,right=6pt,top=3pt,bottom=3pt,before skip=4pt,after skip=2pt,fontupper=\small,
    overlay={\node[anchor=west] at ([xshift=2pt]frame.west) {\includegraphics[width=26pt]{thiam-rond.png}};}]
  #1\end{tcolorbox}}

\newcommand{\C}{\mathbb{C}}
\newcommand{\entete}[3]{% titre, chapitre, classe
\vspace*{-2pt}
\begin{tcolorbox}[enhanced,colback=fond,colframe=ink,boxrule=1pt,arc=3mm,left=6pt,right=6pt,top=4pt,bottom=4pt]
\begin{minipage}[c]{0.15\textwidth}\centering\includegraphics[height=2.25cm]{logo-full.jpg}\end{minipage}%
\begin{minipage}[c]{0.70\textwidth}\centering
  {\small\color{sky}\bfseries COURS DE RENFORCEMENT · ANNÉE SCOLAIRE 2026 -- 2027}\\[3pt]
  {\LARGE\color{ink}\bfseries #1}\\[2pt]
  {\large\color{ink}\bfseries #2}\\[4pt]
  \tikz\node[fill=lime,text=white,rounded corners=2pt,inner sep=3pt,font=\small\bfseries]{Classe~: #3};\
  \tikz\node[fill=sky,text=white,rounded corners=2pt,inner sep=3pt,font=\small\bfseries]{Prof~: M. Racine Thiam};\
  \tikz\node[fill=gold,text=white,rounded corners=2pt,inner sep=3pt,font=\small\bfseries]{CEM Kennedy};
\end{minipage}%
\begin{minipage}[c]{0.15\textwidth}\centering\includegraphics[height=2.25cm]{profdemaths.jpg}\end{minipage}
\end{tcolorbox}}
\newenvironment{competences}{%
\begin{tcolorbox}[enhanced,colback=sky!6,colframe=sky,boxrule=0.7pt,arc=2mm,left=6pt,right=6pt,bottom=3pt,
  title={\bfseries\small Compétences travaillées dans cette fiche},fonttitle=\small,colbacktitle=sky,coltitle=white,
  attach boxed title to top left={xshift=8pt,yshift=-2mm},boxed title style={arc=1mm,boxrule=0pt},top=6pt]
\small\begin{multicols}{2}\begin{itemize}}{%
\end{itemize}\end{multicols}\vspace{-6pt}
{\footnotesize\color{ink!80}Niveau~: \etoiles1\ application directe \quad \etoiles2\ entraînement \quad \etoiles3\ approfondissement \qquad \demo\ exercice de raisonnement}
\end{tcolorbox}}
\newcommand{\rappel}[1]{\begin{tcolorbox}[enhanced,colback=lime!7,colframe=lime!70!black,boxrule=0.6pt,arc=2mm,left=6pt,right=6pt,top=3pt,bottom=3pt,before skip=3pt,after skip=3pt,fontupper=\small,title={\bfseries\footnotesize Rappel de cours},colbacktitle=lime!70!black,coltitle=white,attach boxed title to top left={xshift=8pt,yshift=-2mm},boxed title style={arc=1mm,boxrule=0pt},top=5pt]#1\end{tcolorbox}}

\begin{document}
\entete{FICHE DE TD N°1}{ÉQUATIONS, INÉQUATIONS, SYSTÈMES ET POLYNÔMES}{1ère S2}

\begin{competences}
  \item Résoudre équations et inéquations du \textbf{second degré} et se ramener au second degré.
  \item Résoudre des équations \textbf{irrationnelles} et avec \textbf{valeur absolue}.
  \item Utiliser la \textbf{somme} et le \textbf{produit} des racines, discuter une équation à \textbf{paramètre}.
  \item Résoudre des \textbf{systèmes} linéaires ($2\times2$, pivot de Gauss) et non linéaires.
  \item \textbf{Mathématiser} un problème concret et vérifier la solution.
  \item \textbf{Polynômes}~: racines, factorisation, Horner, identification, division euclidienne.
\end{competences}

% ================= PARTIE A =================
\partie{ink}{A · Équations et inéquations du second degré}{« Discriminant d'abord, conclusion ensuite »}

\begin{exo}{ink}{Échauffement \etoiles1}
Résoudre dans $\R$~:
\[
\text{a) } x^2-6x+5=0 \qquad \text{b) } 4x^2-4x+1=0 \qquad \text{c) } 2x^2+x+3=0 \qquad \text{d) } -x^2+4x+5>0
\]
\[
\text{e) } (3x-2)^2-(x+4)^2\leq0 \qquad\quad \text{f) } \frac{x-1}{x+3}\geq 2 \qquad\quad \text{g) } \frac{x^2-4}{x^2+x-6}<0
\]
\end{exo}

\begin{exo}{ink}{Somme et produit des racines \etoiles2}
\begin{enumerate}
  \item Justifier que l'équation $x^2-5x+3=0$ admet deux solutions $x'$ et $x''$. \textbf{Sans les calculer}, déterminer~:
  \[ x'+x''\ ;\quad x'x''\ ;\quad x'^2+x''^2\ ;\quad \frac1{x'}+\frac1{x''}\ ;\quad (x'-x'')^2\ ;\quad x'^2x''+x'x''^2. \]
  \item Vérifier que $-5$ est solution de $2x^2+7x-15=0$, puis trouver l'autre solution \textbf{sans calculer $\Delta$}.
  \item Un champ rectangulaire a un périmètre de $34$~m et une aire de $60$~m$^2$. Déterminer ses dimensions.
\end{enumerate}
\end{exo}

\begin{exo}{ink}{Changements de variable \etoiles2}
Résoudre dans $\R$~:
\[
\text{a) } x^4-10x^2+9=0 \qquad \text{b) } 2x-7\sqrt x+3=0 \qquad \text{c) } (x^2+x)^2-8(x^2+x)+12=0 \qquad \text{d) } x^4-5x^2+4\leq0
\]
\end{exo}

\begin{exo}{ink}{Radicaux et valeurs absolues \etoiles2\ \demo}
Résoudre dans $\R$ en justifiant les équivalences~:
\[
\text{a) } \sqrt{3x+1}=x-1 \qquad \text{b) } \sqrt{x+2}\leq x \qquad \text{c) } |2x-3|=x+1 \qquad \text{d) } |x^2-5x|=6
\]
\end{exo}
\bulle{\textbf{Attention aux faux amis~!} Pour $\sqrt A=B$, on écrit $\big(B\geq0$ \textbf{et} $A=B^2\big)$. Sans la condition $B\geq0$, tu récupères des solutions qui n'en sont pas. Le bon réflexe~: \textbf{toujours vérifier} tes solutions dans l'équation de départ.}

\begin{exo}{ink}{Inéquations produits et quotients \etoiles2}
Résoudre dans $\R$ à l'aide d'un tableau de signes~:
\[
\text{a) } \frac{(x^2-3x+2)(x^2-9)}{x^2+1}\geq0 \qquad \text{b) } \frac1{x-1}-\frac1{x+1}>1 \qquad \text{c) } \frac{2x+1}{x-2}\leq\frac{x+1}{x+2}
\]
\end{exo}

\begin{exo}{ink}{Problèmes du second degré \etoiles2}
\begin{enumerate}
  \item Trouver un réel qui, ajouté à son inverse, donne $\dfrac{10}{3}$.
  \item Trouver deux entiers consécutifs dont la somme des carrés vaut $313$.
  \item Un jardin carré de Sigma Cap est entouré d'une allée de $1$~m de large. L'aire de l'allée est égale à $44$~m$^2$. Quel est le côté du jardin~?
\end{enumerate}
\end{exo}

\begin{exo}{ink}{Équation à paramètre \etoiles3}
Soit $m$ un réel et \ $(E_m)~:\ x^2-2(m+1)x+m^2+3=0$.
\begin{enumerate}
  \item Montrer que le discriminant réduit de $(E_m)$ est $\Delta'=2m-2$.
  \item Discuter, suivant les valeurs de $m$, l'existence et le \textbf{signe} des solutions de $(E_m)$. Que remarque-t-on~?
  \item Pour quelle valeur de $m$ l'équation a-t-elle une solution double~? La donner.
  \item Déterminer $m$ pour que $3$ soit solution de $(E_m)$.
  \item Trouver une relation indépendante de $m$ entre la somme $S$ et le produit $P$ des solutions.
  \item Déterminer $m$ pour que $x'^2+x''^2=20$ \ (attention aux valeurs à rejeter~!).
  \item Peut-on trouver $m$ pour que les solutions soient \textbf{opposées}~? \textbf{inverses}~? Justifier.
  \item Pour quelles valeurs de $m$ a-t-on \ $x^2-2(m+1)x+m^2+3>0$ \ pour tout réel $x$~?
\end{enumerate}
\end{exo}
\bulle{\textbf{Méthode de discussion~:} 1) signe de $\Delta$~; 2) si $\Delta\geq0$, signe du produit $P$ (solutions de même signe ou non)~; 3) si $P>0$, signe de la somme $S$. Et une valeur de $m$ trouvée par calcul n'est valable que si $\Delta\geq0$~: vérifie toujours~!}

\begin{exo}{ink}{Démonstration~: somme et produit \demo\ \etoiles2}
Soit l'équation $ax^2+bx+c=0$ ($a\neq0$) de discriminant $\Delta\geq0$ et de solutions $x'=\dfrac{-b-\sqrt\Delta}{2a}$ et $x''=\dfrac{-b+\sqrt\Delta}{2a}$.
\begin{enumerate}
  \item Calculer $x'+x''$. Calculer $x'x''$ en utilisant l'identité $(u-v)(u+v)=u^2-v^2$.
  \item En déduire que si $a$ et $c$ sont de signes contraires, l'équation a toujours deux solutions de signes contraires.
\end{enumerate}
\end{exo}

% ================= PARTIE B =================
\partie{lime!80!black}{B · Systèmes d'équations}{« Une inconnue bien choisie, un problème résolu »}

\begin{exo}{lime!80!black}{Systèmes linéaires $2\times2$ \etoiles1}
\begin{enumerate}
  \item Résoudre par la méthode des déterminants~: \ $\begin{cases}2x+3y=12\\3x-y=7\end{cases}$.
  \item Discuter suivant le réel $m$ le système \ $\begin{cases}x+my=1\\mx+4y=2\end{cases}$.
\end{enumerate}
\end{exo}

\begin{exo}{lime!80!black}{Pivot de Gauss et problème \etoiles2}
\begin{enumerate}
  \item Résoudre par la méthode du pivot de Gauss~: \ $\begin{cases}x+y+z=2\\x-y+2z=9\\2x+y-z=-3\end{cases}$.
  \item Au marché de Colobane, Awa achète $2$~kg de mangues, $1$~kg d'oranges et $3$~kg de bananes pour $3\,500$~F. Fatou achète $1$~kg de mangues, $2$~kg d'oranges et $1$~kg de bananes pour $2\,400$~F. Moussa achète $3$~kg de mangues, $2$~kg d'oranges et $2$~kg de bananes pour $4\,300$~F. Déterminer le prix d'un kilogramme de chaque fruit.
\end{enumerate}
\end{exo}

\begin{exo}{lime!80!black}{Systèmes non linéaires \etoiles2}
Résoudre dans $\R^2$~:
\[
\begin{cases}x+y=11\\xy=28\end{cases}\qquad
\begin{cases}x-y=3\\xy=10\end{cases}\qquad
\begin{cases}x^2+y^2=13\\x+y=5\end{cases}
\]
\end{exo}

\begin{exo}{lime!80!black}{Systèmes avec changement d'inconnues \etoiles3}
Résoudre dans $\R^2$ en posant de nouvelles inconnues~:
\[
\begin{cases}\dfrac2x-\dfrac1y=3\\[4pt]\dfrac1x+\dfrac3y=5\end{cases}
\qquad\qquad
\begin{cases}\sqrt x+\sqrt y=5\\x+y=13\end{cases}
\qquad\qquad
\begin{cases}x^2+y^2=20\\xy=8\end{cases}
\]
\end{exo}

\begin{exo}{lime!80!black}{Problème d'âges \etoiles2}
Aujourd'hui, M. Diop a le triple de l'âge de son fils. Dans $12$ ans, il aura le double de l'âge de son fils.
\begin{enumerate}
  \item Choisir les inconnues (une phrase pour chacune, avec l'unité) et compléter un tableau « âge aujourd'hui / âge dans 12 ans ».
  \item Traduire chaque phrase par une équation, résoudre le système et conclure. Vérifier.
\end{enumerate}
\end{exo}
\bulle{\textbf{Mathématiser, c'est traduire~:} « le double de l'âge du fils dans 12 ans » s'écrit $2(y+12)$ et \textbf{pas} $2y+12$~: les parenthèses sont importantes~! Termine toujours par une phrase qui répond à la question.}

% ================= PARTIE C =================
\partie{sky}{C · Polynômes}{« Factoriser, c'est simplifier la vie »}

\begin{exo}{sky}{Développer, réduire, ordonner \etoiles1}
\begin{enumerate}
  \item Développer, réduire et ordonner $P(x)=(x^2-x+1)(x+2)-(x-1)^3$. Quel est son degré~? Est-ce surprenant~?
  \item Déterminer les réels $a$, $b$, $c$ tels que pour tout réel $x$~: $x^3-7x+6=(x-1)(ax^2+bx+c)$, puis factoriser complètement $x^3-7x+6$.
\end{enumerate}
\end{exo}

\begin{exo}{sky}{Racine évidente, Horner \etoiles2}
\begin{enumerate}
  \item Soit $P(x)=x^3-4x^2+x+6$. Vérifier que $-1$ est racine de $P$, factoriser $P$ à l'aide de la méthode de Horner, puis résoudre $P(x)>0$.
  \item Soit $Q(x)=2x^3+x^2-13x+6$. Trouver une racine entière de $Q$, factoriser $Q$ et résoudre $Q(x)\leq0$.
  \item Déterminer le réel $m$ pour que $2$ soit racine de $R(x)=x^3+mx^2-4x+4$, puis factoriser $R(x)$.
\end{enumerate}
\end{exo}

\begin{exo}{sky}{Bicarré et identification \etoiles2}
\begin{enumerate}
  \item Factoriser $P(x)=x^4-13x^2+36$ en produit de quatre polynômes de degré $1$, puis résoudre $P(x)<0$.
  \item Déterminer le polynôme $T$ de degré $2$ tel que $T(0)=3$, $T(1)=0$ et $T(-1)=8$, puis factoriser $T(x)$.
\end{enumerate}
\end{exo}

\begin{exo}{sky}{Division euclidienne \etoiles2}
Effectuer la division euclidienne de \ $A(x)=2x^3-3x^2+4x-5$ \ par \ $B(x)=x^2-x+1$. Écrire l'égalité obtenue et la vérifier pour $x=0$ et $x=1$.
\end{exo}

\begin{exo}{sky}{Pourquoi une racine donne un facteur \demo\ \etoiles3}
\begin{enumerate}
  \item Démontrer que pour tout entier $k\geq1$ et tous réels $x$ et $a$~: \ $x^k-a^k=(x-a)\left(x^{k-1}+x^{k-2}a+\cdots+xa^{k-2}+a^{k-1}\right)$.
  \item Soit $P(x)=c_nx^n+\cdots+c_1x+c_0$ et $a$ une racine de $P$. En écrivant $P(x)=P(x)-P(a)$, démontrer qu'il existe un polynôme $Q$ tel que $P(x)=(x-a)Q(x)$.
  \item \textbf{Application~:} déterminer tous les polynômes de degré au plus $2$ qui s'annulent en $-2$ et en $1$. Quel est celui qui prend la valeur $6$ en $0$~?
\end{enumerate}
\end{exo}

\begin{exo}{sky}{Polynôme réciproque de degré 3 \etoiles3}
Soit $P(x)=2x^3+3x^2+3x+2$.
\begin{enumerate}
  \item Remarquer la symétrie des coefficients. Vérifier que $-1$ est racine de $P$.
  \item Factoriser $P(x)$ et résoudre $P(x)=0$, puis $P(x)<0$.
\end{enumerate}
\end{exo}
\bulle{\textbf{Astuce~:} quand les coefficients sont symétriques ($2,3,3,2$) et le degré impair, $-1$ est \textbf{toujours} racine. Pense à tester $1$, $-1$, $2$, $-2$… avant de te lancer dans des calculs compliqués~!}

% ================= PARTIE D =================
\partie{violet}{D · Fractions rationnelles et synthèse}{« On rassemble tout »}

\begin{exo}{violet}{Fractions rationnelles \etoiles2\ \demo}
\begin{enumerate}
  \item Déterminer les réels $a$, $b$, $c$ tels que pour tout $x\neq-1$~: \ $\dfrac{x^2+3x+4}{x+1}=ax+b+\dfrac{c}{x+1}$.
  \item \demo\ Démontrer que pour tout entier $k\geq1$~: $\dfrac1{k(k+1)}=\dfrac1k-\dfrac1{k+1}$. En déduire, pour $n\geq1$, la somme
  \[ S_n=\frac1{1\times2}+\frac1{2\times3}+\cdots+\frac1{n(n+1)}=\frac{n}{n+1}. \]
\end{enumerate}
\end{exo}

\begin{exo}{violet}{Une racine qui ne bouge pas \etoiles3}
Soit $m$ un réel et \ $(E_m)~:\ (m-2)x^2-2(m-1)x+m=0$.
\begin{enumerate}
  \item Résoudre $(E_m)$ pour $m=2$.
  \item Pour $m\neq2$, vérifier que $1$ est solution de $(E_m)$ quel que soit $m$. À l'aide du produit des solutions, en déduire l'autre solution $x''$.
  \item Calculer le discriminant réduit $\Delta'$ et confirmer le résultat. L'équation peut-elle avoir une solution double~?
  \item Déterminer les valeurs de $m$ pour lesquelles $x''>1$, puis celles pour lesquelles $x''$ est négative.
\end{enumerate}
\end{exo}

\begin{exo}{violet}{Vrai ou faux~? Justifier \demo\ \etoiles2}
Pour chaque affirmation, répondre par vrai ou faux et \textbf{justifier} (démonstration ou contre-exemple).
\begin{enumerate}
  \item Si une équation du second degré a un discriminant positif, ses deux solutions sont positives.
  \item Si $a$ et $c$ sont de signes contraires, l'équation $ax^2+bx+c=0$ a deux solutions distinctes.
  \item Le polynôme $x^3+x+1$ admet une racine entière.
  \item Le polynôme $x^3-3x+2$ est divisible par $(x-1)^2$.
  \item Le système $\begin{cases}x+2y=3\\2x+4y=5\end{cases}$ admet une infinité de solutions.
\end{enumerate}
\end{exo}
\bulle{\textbf{Règle d'or~:} un \textbf{contre-exemple} suffit pour montrer qu'une affirmation est fausse~; pour montrer qu'elle est vraie, il faut une \textbf{démonstration} valable dans tous les cas.}

\begin{exo}{violet}{Problème de synthèse \etoiles3}
Soit \ $P(x)=x^3-3x^2-4x+12$.
\begin{enumerate}
  \item Calculer $P(3)$, puis factoriser complètement $P(x)$.
  \item Soit $f(x)=\dfrac{P(x)}{x^2-5x+6}$. Déterminer l'ensemble de définition $D_f$ de $f$ et simplifier $f(x)$ sur $D_f$.
  \item Résoudre dans $\R$~: \ $f(x)\leq1$, \ puis \ $P(x)\geq0$.
  \item Résoudre le système \ $\begin{cases}x+y=1\\xy=-6\end{cases}$ \ et faire le lien avec les racines de $P$.
\end{enumerate}
\end{exo}

\bulle{\textbf{Le mot de la fin~:} en 1ère S2, la régularité paie~! Refais les exercices où tu as bloqué, vérifie tes solutions, et rédige avec des phrases (« on sait que… donc… »). La réussite, c'est maintenant~! \hfill\textbf{M. Thiam}}

\end{document}
